% Terjemahan Bahasa Indonesia dari chapter1.tex baris 568--629.
% Sumber: Methods of Algebra, Volume 2, commit
% 9a5803ff77dd3257484cb177f851a73770a59dd3 (CC BY 4.0).
% stable-unit-id: o014.aljabr2.chapter1.linear-categories-over-commutative-rings

% segment-id: o014.li2.chapter1.linear-categories-over-commutative-rings.h001
\section{Perumuman: Kategori Linear atas Gelanggang Komutatif}\label{sec:linear-cat}
\hypertarget{o014.aljabr2.chapter1.linear-categories-over-commutative-rings}{ }

% segment-id: o014.li2.chapter1.linear-categories-over-commutative-rings.p001
Himpunan-himpunan $\Hom$ dalam suatu kategori aditif mempunyai struktur grup
abelian. Dalam banyak situasi, himpunan-himpunan $\Hom$ ini juga dilengkapi
dengan perkalian skalar oleh suatu gelanggang komutatif $\Bbbk$. Sebagai
contoh, dalam $\Bbbk\dcate{Mod}$ setiap himpunan $\Hom$ secara alami merupakan
modul-$\Bbbk$, dan komposisi morfisme bersifat $\Bbbk$-bilinear.

% segment-id: o014.li2.chapter1.linear-categories-over-commutative-rings.d001
\begin{definisi}\label{def:cat-k-linear}
	\index{kategori!$\Bbbk\dcate{Mod}$-}
	\index{linear!$\Bbbk$- ($\Bbbk$-linear)}
	Misalkan $\Bbbk$ merupakan gelanggang komutatif. Andaikan setiap himpunan
	$\Hom$ dalam kategori $\mathcal{A}$ dilengkapi dengan struktur
	modul-$\Bbbk$ sedemikian sehingga komposisi morfisme
	$\Hom(Y, Z) \times \Hom(X, Y) \to \Hom(X, Z)$ merupakan pemetaan
	$\Bbbk$-bilinear untuk semua objek $X,Y,Z$; dengan kata lain, diagram
	berikut komutatif:
	% segment-id: o014.li2.chapter1.linear-categories-over-commutative-rings.g001
	\[\begin{tikzcd}
		\Hom(Y, Z) \times \Hom(X, Y) \arrow[d] \arrow[r, "\text{komposisi}"] & \Hom(X, Z) \\
		\Hom(Y, Z) \dotimes{\Bbbk} \Hom(X, Y) \arrow[ru, bend right, start anchor=east, "{\exists! \; \text{homomorfisme modul-$\Bbbk$}}"']
	\end{tikzcd}  \]
	Dengan data ini, $\mathcal{A}$ disebut \emph{kategori-$\Bbbk\dcate{Mod}$}.

	Misalkan $F: \mathcal{A} \to \mathcal{A}'$ merupakan funktor antara
	kategori-$\Bbbk\dcate{Mod}$. Jika pemetaan
	$\Hom(X,Y) \to \Hom(FX,FY)$ merupakan homomorfisme modul-$\Bbbk$ untuk
	semua $X,Y \in \Obj(\mathcal{A})$, maka $F$ disebut
	\emph{funktor $\Bbbk$-linear}.
\end{definisi}

% segment-id: o014.li2.chapter1.linear-categories-over-commutative-rings.p002
Perhatikan bahwa $\Bbbk\dcate{Mod}$ merupakan kategori monoidal terhadap hasil
kali tensor $\otimes := \dotimes{\Bbbk}$. Oleh karena itu, istilah
``kategori-$\Bbbk\dcate{Mod}$'' konsisten dengan terminologi pengayaan kategori
dalam \cite[\S 3.4]{Li1}. Kategori-$\cate{Ab}$ yang diperkenalkan sebelumnya
tepat merupakan kategori-$\Z\dcate{Mod}$. Sebaliknya, dengan melupakan
perkalian skalar pada suatu kategori-$\Bbbk\dcate{Mod}$, kita memperoleh
kategori-$\cate{Ab}$.

% segment-id: o014.li2.chapter1.linear-categories-over-commutative-rings.e001
\begin{contoh}
	Kategori-$\Bbbk\dcate{Mod}$ yang hanya mempunyai satu objek tepat
	merupakan aljabar-$\Bbbk$. Korespondensinya diberikan oleh
	$\mathcal{A} \mapsto \End_{\mathcal{A}}(\star)$, dengan
	$\Obj(\mathcal{A}) = \{\star\}$.
\end{contoh}

% segment-id: o014.li2.chapter1.linear-categories-over-commutative-rings.p003
Argumen dalam Proposisi~\ref{prop:automatic-Ab-morphism} dapat diterapkan tanpa
perubahan dan memberikan hasil berikut.

% segment-id: o014.li2.chapter1.linear-categories-over-commutative-rings.pr001
\begin{proposisi}\label{prop:automatic-k-morphism}
	Misalkan $\mathcal{A}$ dan $\mathcal{A}'$ merupakan
	kategori-$\Bbbk\dcate{Mod}$.
	% segment-id: o014.li2.chapter1.linear-categories-over-commutative-rings.l001
	\begin{enumerate}
		% segment-id: o014.li2.chapter1.linear-categories-over-commutative-rings.i001
		\item Diberikan sepasang funktor adjoin
		% Diagram adjungsi sumber yang semula sebaris direflow sebagai display
		% terpusat agar tidak menyempitkan atau melampaui lebar baris pembaca.
		% Rantai kesetaraan panjang dalam bukti juga direflow sebagai display.
		% segment-id: o014.li2.chapter1.linear-categories-over-commutative-rings.g002
		\[
		\begin{tikzcd}
			F: \mathcal{A} \arrow[shift left, r] & \mathcal{A}' \arrow[shift left, l] : G
		\end{tikzcd},
		\]
		dengan $F$ dan $G$ keduanya bersifat $\Bbbk$-linear. Maka isomorfisme adjungsi
		$\Hom_{\mathcal{A}'}(F(\cdot),\cdot) \rightiso
		\Hom_{\mathcal{A}}(\cdot,G(\cdot))$ bersifat $\Bbbk$-linear.
		% segment-id: o014.li2.chapter1.linear-categories-over-commutative-rings.i002
		\item Untuk setiap $\alpha: I \to \mathcal{A}$, bijeksi
		% segment-id: o014.li2.chapter1.linear-categories-over-commutative-rings.q001
		\[
			\Hom_{\mathcal{A}}\left(\varinjlim \alpha,T\right)
			\simeq \varprojlim_{i \in \Obj(I)}
			\Hom_{\mathcal{A}}(\alpha(i),T),
			\quad T \in \Obj(\mathcal{A})
		\]
		bersifat $\Bbbk$-linear, asalkan $\varinjlim \alpha$ ada. Pernyataan yang sama
		berlaku untuk $\varprojlim$.
	\end{enumerate}
\end{proposisi}

% segment-id: o014.li2.chapter1.linear-categories-over-commutative-rings.p004
Kecuali dinyatakan lain, mulai sekarang semua funktor dan ekuivalensi antara
kategori-$\Bbbk\dcate{Mod}$ dipahami bersifat $\Bbbk$-linear.

% segment-id: o014.li2.chapter1.linear-categories-over-commutative-rings.d002
\begin{definisi}\label{def:k-linear-cat}
	\index{kategori!$\Bbbk$-linear ($\Bbbk$-linear)}
	Jika suatu kategori aditif $\mathcal{A}$ dilengkapi dengan struktur
	$\Bbbk$-linear yang kompatibel dengan struktur kategori-$\cate{Ab}$ yang
	sudah ada, maka $\mathcal{A}$ beserta data tersebut disebut
	\emph{kategori $\Bbbk$-linear}.
\end{definisi}

% segment-id: o014.li2.chapter1.linear-categories-over-commutative-rings.p005
Hal ini sama artinya dengan mengganti kategori-$\cate{Ab}$ dengan
kategori-$\Bbbk\dcate{Mod}$ dalam pembahasan \S\ref{sec:Ker-Coker}. Sebagian
besar hasil mengenai kategori aditif dalam bab ini dapat diperluas ke situasi
$\Bbbk$-linear dengan argumen yang sama tanpa perubahan. Satu-satunya pengecualian
berkaitan dengan Korolari~\ref{prop:automatic-additivity}: dalam kategori
aditif, penjumlahan pada himpunan $\Hom$ dapat dicirikan oleh sifat kategori itu
sendiri, tetapi hal ini tidak berlaku bagi perkalian skalar oleh $\Bbbk$.
Karena itu, tidak ada alasan bagi funktor (bahkan ekuivalensi) antara kategori
$\Bbbk$-linear untuk otomatis bersifat $\Bbbk$-linear.

% segment-id: o014.li2.chapter1.linear-categories-over-commutative-rings.p006
Berikut ini beberapa hasil mengenai sifat $\Bbbk$-linear. Ingat bahwa
\emph{pusat} suatu kategori $\mathcal{A}$ didefinisikan sebagai
\index{pusat (center)} \index[sym1]{ZA@$Z(\mathcal{A})$}
$Z(\mathcal{A}) := \End(\identity_{\mathcal{A}})$. Unsur-unsurnya merupakan
morfisme dari funktor identitas $\identity_{\mathcal{A}}$ ke dirinya sendiri,
yang ditulis sebagai $(a_X)_{X \in \Obj(\mathcal{A})}$ dengan
$a_X \in \End_{\mathcal{A}}(X)$. Dengan operasi komposisi morfisme, pusat tersebut
merupakan monoid komutatif; lihat \cite[Definisi 2.3.8,
Proposisi 2.3.9]{Li1}. Jika $\mathcal{A}$ merupakan kategori-$\cate{Ab}$, maka
$Z(\mathcal{A})$ menjadi gelanggang komutatif terhadap komposisi dan
penjumlahan
$(a_X)_X + (b_X)_X := (a_X+b_X)_X$.

% segment-id: o014.li2.chapter1.linear-categories-over-commutative-rings.pr002
\begin{proposisi}\label{prop:k-linear-center}
	Misalkan $\mathcal{A}$ merupakan kategori-$\cate{Ab}$. Maka terdapat
	bijeksi
	% segment-id: o014.li2.chapter1.linear-categories-over-commutative-rings.q002
	\[
		\left\{\text{struktur kategori-$\Bbbk\dcate{Mod}$ pada $\mathcal{A}$}\right\}
		\stackrel{1:1}{\longleftrightarrow}
		\Hom_{\text{gelanggang}}(\Bbbk,Z(\mathcal{A})).
	\]
	Pemetaan tersebut diberikan sebagai berikut. Jika struktur
	kategori-$\Bbbk\dcate{Mod}$ pada $\mathcal{A}$ telah ditetapkan, untuk setiap
	$a \in \Bbbk$ dan $X \in \Obj(\mathcal{A})$ definisikan
	$a_X := a \cdot \identity_X \in \End_{\mathcal{A}}(X)$. Maka
	$a \mapsto (a_X)_{X \in \Obj(\mathcal{A})} \in Z(\mathcal{A})$ merupakan
	homomorfisme gelanggang. Sebaliknya, diberikan homomorfisme gelanggang di
	ruas kanan,
	$a \mapsto (a_X)_{X \in \Obj(\mathcal{A})}$, definisikan struktur
	kategori-$\Bbbk\dcate{Mod}$ pada $\mathcal{A}$ dengan
	$a \cdot f := f \circ a_X = a_Y \circ f$, untuk
	$f \in \Hom(X,Y)$.
\end{proposisi}

% segment-id: o014.li2.chapter1.linear-categories-over-commutative-rings.pf001
\begin{bukti}
	Syarat yang mencirikan
	$(a_X)_{X \in \Obj(\mathcal{A})} \in Z(\mathcal{A})$ adalah
	$f \circ a_X = a_Y \circ f$ untuk setiap $f \in \Hom(X,Y)$. Jika
	$\mathcal{A}$ dilengkapi dengan struktur kategori-$\Bbbk\dcate{Mod}$,
	definisikan $a_X := a \cdot \identity_X$. Sifat bilinear memberikan
	\[
		f \circ a_X = f \circ (a \cdot \identity_X)
		= (a f) \circ \identity_X
		= (a \cdot \identity_Y) \circ f
		= a_Y \circ f,
	\]
	sehingga $(a_X)_{X \in \Obj(\mathcal{A})} \in Z(\mathcal{A})$.
	Verifikasi lainnya langsung.
\end{bukti}

% segment-id: o014.li2.chapter1.linear-categories-over-commutative-rings.p007
Sebagai akibatnya, kategori-$\cate{Ab}$ $\mathcal{A}$ secara alami menjadi
kategori-$Z(\mathcal{A})\dcate{Mod}$.
